\section{Further research questions with concrete acceptance criteria}
\label{sec:further}

\paragraph{1. Can a source invariant pay for one actual binary state?}
A forced-zero decrement is valuable only if replacing its control
branch changes a compiled runtime component enough to delete a
working state.  The acceptance criterion is a rebuilt complete
table with at most 236 states, a checked source transformation,
new decoder coverage and initialization, and either a direct
unbounded runtime replay or an applicable representation theorem.
A smaller leaf list alone is insufficient.

\paragraph{2. Can two state changes be synthesized together?}
The 237-state fold succeeds because a retained helper's return
changes along with the state identification.  A useful search
should consider a small connected group of controls, permit
several guarded representations, and verify all incoming
contexts together.  A SAT/SMT encoding could separate the
finite candidate synthesis from the independent positive-block
checker.  The checker, not the solver's empirical profile,
must discharge the all-tape or valid-frame theorem actually used.

\paragraph{3. Can a stronger source abstraction preserve useful correlations?}
Independent register sets lose relations such as equality,
one-step offsets, and a scratch register being zero whenever
another control branch is taken.  An economical relational
domain might expose additional dead code or simpler arithmetic
helpers.  Its transitions, join operation and initial inclusion
must have an ordinary soundness proof and an independent finite
closure check before those facts change the proof-search program.

\paragraph{4. What is the fully charged cost of E692?}
Implement one fair ground-rewrite verifier, including the union
of variables of both equation sides, contextual replacement,
enumeration, reset and a writer for the complete 692-bit payload.  Compare the final
binary state count with the existing arithmetic engine.
An evaluator with a small code file but a large tape loader
does not meet the target.  The most useful intermediate output
would be a small fixed register program with exact source
semantics and a fully included packet.

\paragraph{5. Can the strongest ordinary proofs be given compact kernel endpoints?}
The final three-block local fold is an especially small target:
its invariant is elementary and every exceptional path has at
most five steps.  The equational construction has a different
target: formalize finite descriptor descent, the weighted raw
assignment, and the complemented-witness lemma in one development.
These tasks would strengthen reproducibility without implying
that the unresolved internal PA proof-code comparison has
already been completed.

The principal open outcome remains the original one: an explicit
sub-237-state blank binary machine with a complete calibration
argument.  The validated checkpoint and the new finite
consistency problem provide concrete starting points for that work.
